% GENERATED FILE. Edit canonical lessons, then run npm run generate:books.
% canonical-source-hash: fnv1a64:40b307641e26fbcd
% canonical-lessons: AR-C534-ataa, AR-C534-asa, AR-C534-wabbakha, AR-C534-ihtiram, AR-C534-afa

\chapter{Obeying and forgiving}
\label{ch:ar-obeying-and-forgiving}
\hlchaptermodality{534}

\begin{chapteropening}
\textbf{By the end of this chapter:} \emph{I can say that someone obeyed, disobeyed, scolded or pardoned someone, and talk about respect.}

Complete the last lesson of chapter 534: I can say that someone obeyed, disobeyed, scolded or pardoned someone, and talk about respect.
\end{chapteropening}

\section[\texorpdfstring{\ar{أطاع} (aṭāʿa)}{أطاع (aṭāʿa)}]{\ar{أطاع} (aṭāʿa) — to obey}

\label{lesson:AR-C534-ataa}



\begin{quote}
Before the new one: say the Arabic for to stand by, then the Arabic for a loved one.
\end{quote}

\subsection*{You'll want to know: \ar{أطاع}}
\textbf{\ar{أطاع}} — \emph{aṭāʿa} — \textquotedblleft{}to obey\textquotedblright{}.

\begin{quote}
\ar{أطاع} \ar{الولد} \ar{أباه} \ar{وذهب} \ar{إلى} \ar{النوم}.
\end{quote}

\begin{quote}
\emph{aṭāʿa al-walad abāhu wa-dhahaba ilā an-nawm.}
\end{quote}

\begin{quote}
The boy obeyed his father and went to bed.
\end{quote}

Present: \ar{يطيع}, \emph{yuṭīʿ}.

\begin{cousinweb}
\textbf{\ar{أطاع}} (\emph{aṭāʿa}) is built on the native Arabic root \emph{ṭ-w-ʿ}, \textquotedblleft{}to comply, to be willing\textquotedblright{} (Lane). Istaṭāʿa, to be able, which you know, is from the same root.
\end{cousinweb}

\begin{grammarlens}[title={What you've built}]
One more word for describing what you lived and felt.
\end{grammarlens}

\subsection*{Your turn}
\begin{itemize}
\raggedright
  \item \emph{Say it:} \emph{aṭāʿa}
  \item \emph{Say it:} \emph{aṭāʿa}, once more
  \item \emph{Say it:} \emph{aṭāʿa al-walad abāhu wa-dhahaba ilā an-nawm.} — The boy obeyed his father and went to bed.
  \item \emph{Recall:} say the Arabic for to stand by, then the Arabic for a loved one, then say \emph{aṭāʿa} again
\end{itemize}

\begin{tcolorbox}[breakable,colback=teal!4,colframe=teal!35!black,title={Before you move on}]
What does \ar{أطاع} mean? (\textquotedblleft{}To obey\textquotedblright{}.) Say it once more.
\end{tcolorbox}
\section[\texorpdfstring{\ar{عصى} (ʿaṣā)}{عصى (ʿaṣā)}]{\ar{عصى} (ʿaṣā) — to disobey}

\label{lesson:AR-C534-asa}



\begin{quote}
Before the new one: say the Arabic for a loved one, then the Arabic for to obey.
\end{quote}

\subsection*{You'll want to know: \ar{عصى}}
\textbf{\ar{عصى}} — \emph{ʿaṣā} — \textquotedblleft{}to disobey\textquotedblright{}.

\begin{quote}
\ar{عصت} \ar{البنت} \ar{أمها} \ar{وخرجت} \ar{في} \ar{المطر}.
\end{quote}

\begin{quote}
\emph{ʿaṣat al-bint ummahā wa-kharajat fī al-maṭar.}
\end{quote}

\begin{quote}
The girl disobeyed her mother and went out in the rain.
\end{quote}

Present: \ar{يعصي}, \emph{yaʿṣī}.

\begin{cousinweb}
\textbf{\ar{عصى}} (\emph{ʿaṣā}) is built on the native Arabic root \emph{ʿ-ṣ-y}, \textquotedblleft{}to resist, to rebel\textquotedblright{} (Lane). Its last root letter y shows as alif in the past and as ī in the present.
\end{cousinweb}

\begin{grammarlens}[title={What you've built}]
One more word for describing what you lived and felt.
\end{grammarlens}

\subsection*{Your turn}
\begin{itemize}
\raggedright
  \item \emph{Say it:} \emph{ʿaṣā}
  \item \emph{Say it:} \emph{ʿaṣā}, once more
  \item \emph{Say it:} \emph{ʿaṣat al-bint ummahā wa-kharajat fī al-maṭar.} — The girl disobeyed her mother and went out in the rain.
  \item \emph{Recall:} say the Arabic for a loved one, then the Arabic for to obey, then say \emph{ʿaṣā} again
\end{itemize}

\begin{tcolorbox}[breakable,colback=teal!4,colframe=teal!35!black,title={Before you move on}]
What does \ar{عصى} mean? (\textquotedblleft{}To disobey\textquotedblright{}.) Say it once more.
\end{tcolorbox}
\section[\texorpdfstring{\ar{وبخ} (wabbakha)}{وبخ (wabbakha)}]{\ar{وبخ} (wabbakha) — to scold}

\label{lesson:AR-C534-wabbakha}



\begin{quote}
Before the new one: say the Arabic for to obey, then the Arabic for to disobey.
\end{quote}

\subsection*{You'll want to know: \ar{وبخ}}
\textbf{\ar{وبخ}} — \emph{wabbakha} — \textquotedblleft{}to scold\textquotedblright{}.

\begin{quote}
\ar{وبخ} \ar{المدرب} \ar{اللاعبين} \ar{بعد} \ar{الخسارة}.
\end{quote}

\begin{quote}
\emph{wabbakha al-mudarrib al-lāʿibīn baʿda al-khasāra.}
\end{quote}

\begin{quote}
The coach scolded the players after the defeat.
\end{quote}

Present: \ar{يوبخ}, \emph{yuwabbikh}.

\begin{cousinweb}
\textbf{\ar{وبخ}} (\emph{wabbakha}) is built on the native Arabic root \emph{w-b-kh}, \textquotedblleft{}to blame, to rebuke\textquotedblright{} (Lane). Form II wabbakha is to rebuke sharply.
\end{cousinweb}

\begin{grammarlens}[title={What you've built}]
One more word for describing what you lived and felt.
\end{grammarlens}

\subsection*{Your turn}
\begin{itemize}
\raggedright
  \item \emph{Say it:} \emph{wabbakha}
  \item \emph{Say it:} \emph{wabbakha}, once more
  \item \emph{Say it:} \emph{wabbakha al-mudarrib al-lāʿibīn baʿda al-khasāra.} — The coach scolded the players after the defeat.
  \item \emph{Recall:} say the Arabic for to obey, then the Arabic for to disobey, then say \emph{wabbakha} again
\end{itemize}

\begin{tcolorbox}[breakable,colback=teal!4,colframe=teal!35!black,title={Before you move on}]
What does \ar{وبخ} mean? (\textquotedblleft{}To scold\textquotedblright{}.) Say it once more.
\end{tcolorbox}
\section[\texorpdfstring{\ar{احترام} (iḥtirām)}{احترام (iḥtirām)}]{\ar{احترام} (iḥtirām) — esteem, respect}

\label{lesson:AR-C534-ihtiram}



\begin{quote}
Before the new one: say the Arabic for to disobey, then the Arabic for to scold.
\end{quote}

\subsection*{You'll want to know: \ar{احترام}}
\textbf{\ar{احترام}} — \emph{iḥtirām} — \textquotedblleft{}esteem, respect\textquotedblright{}.

\begin{quote}
\ar{يتكلم} \ar{الطلاب} \ar{مع} \ar{المعلم} \ar{باحترام}.
\end{quote}

\begin{quote}
\emph{yatakallamu aṭ-ṭullāb maʿa al-muʿallim bi-iḥtirām.}
\end{quote}

\begin{quote}
The students speak to the teacher with respect.
\end{quote}

The verb \ar{احترم}, to respect, goes with it.

\begin{cousinweb}
\textbf{\ar{احترام}} (\emph{iḥtirām}) is built on the native Arabic root \emph{ḥ-r-m}, \textquotedblleft{}to be forbidden, to be sacred\textquotedblright{} (Lane). Iḥtirām is the verbal noun of Form VIII: treating someone as not to be touched.
\end{cousinweb}

\begin{grammarlens}[title={What you've built}]
One more word for describing what you lived and felt.
\end{grammarlens}

\subsection*{Your turn}
\begin{itemize}
\raggedright
  \item \emph{Say it:} \emph{iḥtirām}
  \item \emph{Say it:} \emph{iḥtirām}, once more
  \item \emph{Say it:} \emph{yatakallamu aṭ-ṭullāb maʿa al-muʿallim bi-iḥtirām.} — The students speak to the teacher with respect.
  \item \emph{Recall:} say the Arabic for to disobey, then the Arabic for to scold, then say \emph{iḥtirām} again
\end{itemize}

\begin{tcolorbox}[breakable,colback=teal!4,colframe=teal!35!black,title={Before you move on}]
What does \ar{احترام} mean? (\textquotedblleft{}Esteem, respect\textquotedblright{}.) Say it once more.
\end{tcolorbox}
\section[\texorpdfstring{\ar{عفا} (ʿafā)}{عفا (ʿafā)}]{\ar{عفا} (ʿafā) — to pardon}

\label{lesson:AR-C534-afa}



\begin{quote}
Before the new one: say the Arabic for to scold, then the Arabic for esteem.
\end{quote}

\subsection*{You'll want to know: \ar{عفا}}
\textbf{\ar{عفا}} — \emph{ʿafā} — \textquotedblleft{}to pardon\textquotedblright{}.

\begin{quote}
\ar{عفا} \ar{المعلم} \ar{عن} \ar{الطالب}.
\end{quote}

\begin{quote}
\emph{ʿafā al-muʿallim ʿan aṭ-ṭālib.}
\end{quote}

\begin{quote}
The teacher pardoned the student.
\end{quote}

Present: \ar{يعفو}, \emph{yaʿfū}.

It takes \ar{عن}.

\begin{cousinweb}
\textbf{\ar{عفا}} (\emph{ʿafā}) is built on the native Arabic root \emph{ʿ-f-w}, \textquotedblleft{}to efface; to pardon; to give health\textquotedblright{} (Lane). ʿAfwan, which you know, is from the same root.
\end{cousinweb}

\begin{grammarlens}[title={What you've built}]
One more word for describing what you lived and felt.
\end{grammarlens}

\subsection*{Your turn}
\begin{itemize}
\raggedright
  \item \emph{Say it:} \emph{ʿafā}
  \item \emph{Say it:} \emph{ʿafā}, once more
  \item \emph{Say it:} \emph{ʿafā al-muʿallim ʿan aṭ-ṭālib.} — The teacher pardoned the student.
  \item \emph{Recall:} say the Arabic for to scold, then the Arabic for esteem, then say \emph{ʿafā} again
\end{itemize}

\begin{tcolorbox}[breakable,colback=teal!4,colframe=teal!35!black,title={Before you move on}]
What does \ar{عفا} mean? (\textquotedblleft{}To pardon\textquotedblright{}.) Say it once more.
\end{tcolorbox}
